\section{The certificates}\label{sec:cert}

\subsection{What is certified}
Four finite statements enter the proofs. Three are local inequalities for all gap vectors in an orthant, one is the majorant.
\begin{itemize}
\item \texttt{CertS8}: $(\mathrm{LI})$ of Definition~\ref{def:LI} for $k_{\tilde\psi}$, $K=8$, with position weights and penalties summing to
$\nu=\frac7{1700}$, at $c=\frac{199}{25000}$ (Theorem~\ref{thm:S}, Corollary~\ref{cor:SC}).
\item \texttt{CertAM5} and \texttt{CertAM7}: $(\mathrm{LI}_{\mathrm m})$ of Definition~\ref{def:LIm} for $k_{\cos1.6s}$ with the claims of
Table~\ref{tab:allmarks} for $K=5$ and $K=7$ (Theorem~\ref{thm:main}); a $K=6$ certificate, shown in Table~\ref{tab:allmarks} for comparison, is not used.
\item \texttt{CertMajV2}: the majorant of \S\ref{sec:simple}: $P\ge0$ on $[0,30]$ and $\cB\ge v_\psi$ on $[0,\frac12]$, with $2\beta_0+4\beta_*<\frac{33}{10}$.
(The library also contains a first version, \texttt{CertMaj}, with the threshold $3.2063638853$, only $5.7\cdot10^{-11}$ above the
actual value; it is proved as well, and no theorem uses it.)
\end{itemize}
The data of the local inequalities (weights, penalties, credits) were chosen by linear programmes over finite sets of gap vectors found by
local minimisation, and the claims were then set slightly below the programme's values: for $K=5$ each claim lies $1.5\cdot10^{-5}$ below the programme's value,
and for $K=7$ every claim was lowered by $10^{-5}$ after the certifier had refuted the programme's own claims (\S\ref{ssec:bnb}). The
certificates, not the programmes, carry the proof. In the formalisation each is a single \texttt{Prop} (Appendix~\ref{app:lean}). The local ones are existential in the weights: they assert
\emph{some} data with the stated $\nu$ and credit sums $a_1,a_2$, since the constants depend on the data only through these numbers; the files
of Table~\ref{tab:certs} supply the witnesses. Gap vectors range over the whole orthant $[0,\infty)^{K-1}$; the reductions used by the
certifiers (to a box, and to classes of patterns) are theorems, not part of the statements.

\begin{table}[ht]\centering\small
\caption{Certificates, their data and their verification. Paths are relative to the folder
\texttt{supplementary/certificates/} of the repository (\S\ref{ssec:reproduce}). The seven sha256 values printed here are those of the
files in \texttt{supplementary/}; those of the three local certificates are also recorded in the doc-strings of their \texttt{Prop}s
(Appendix~\ref{app:defs}).}\label{tab:certs}
\footnotesize
\begin{tabular}{@{}p{1.3cm}>{\raggedright\arraybackslash}p{7.9cm}>{\raggedright\arraybackslash}p{5.8cm}@{}}
\toprule
name & data and sha256 & verification\\
\midrule
\texttt{CertS8} & \texttt{CertS8/\allowbreak w\_poly8A\_K8\_mu1700.json}\newline \texttt{257b340d9da4f343600a21589f16d891}\newline\texttt{5f3e082499d475e153c5453980671e78}\newline the window, \texttt{win\_poly8A.json} of the same folder (the unnormalised $v$):\newline \texttt{c4005f57a7920f78aa2df33b093c9c50}\newline\texttt{30efac2b8e39fc6cf7371b0c0f32dc14} &
interval branch-and-bound with LB3, about $4.5\cdot10^8$ boxes in $8$ shards; re-run from scratch ($4.51\cdot10^8$ boxes). In Lean: displayed hypothesis.\\[4pt]
\texttt{CertAM5} & \texttt{CertAM5/\allowbreak claims\_d15e6\_a1.6\_K5\_mu500.json}\newline \texttt{cf055dd2ee9a9e469b6510fe70b5fbc4}\newline\texttt{16b8a40515cc6746bc390d40d88777f2}\newline the $20$ files \texttt{xpat\_K5\_*.json} of the same folder, concatenated in sorted order:\newline \texttt{a5947c01d34ddcbc7366861b35e83942}\newline\texttt{099fcbf7ac71e13fcd432796882dd34e} &
$20$ classes; $1.18\cdot10^7$ boxes (LB1--LB2), $3.7\cdot10^5$ (LB3); re-run with identical counts. In Lean: replayed by the kernel ($46{,}047$ leaves, $89{,}362$ point records; \S\ref{ssec:replay}), so the $K=5$ pair has no hypothesis; the replay files are the library \texttt{ZetaSReplay} of the repository.\\[4pt]
\texttt{CertAM7} & \texttt{CertAM7/\allowbreak claims\_a1.6\_K7\_mu500\_L1e5.json}\newline \texttt{d2cf393019407bdc797c73147ec8e3dc}\newline\texttt{66ec41fe745600c84173f74ac66bfc5c}\newline the $72$ files \texttt{xpat\_K7\_*.json} of the same folder, concatenated in sorted order:\newline \texttt{33c4c8c5762bcebdf8f544ab34a137a6}\newline\texttt{a7ba4065d08f5caee0a51188e1a67514} &
$72$ classes covering $128$ patterns; about $4.8\cdot10^8$ boxes (LB3); complete run $72/72$ at these claims:
\texttt{logs/cert\_lb3\_rerun.txt} of the same folder, completeness checked in \texttt{logs/k7\_verify\_pexact2.txt}
(the log of the refutation described in \S\ref{ssec:bnb} is not part of \texttt{supplementary/}). In Lean: displayed hypothesis.\\[4pt]
\texttt{CertMajV2} & \texttt{CertMajV2/\allowbreak majorant\_delta1\_beta.npy} ($60$ binary64 numbers)\newline \texttt{fd909f79669e2fc1ba5fd4e44d01d600}\newline\texttt{7ccf5303b7ce62d4c5ee08f7ef024845} &
Arb at $106$ bits. In Lean: \emph{kernel-checked}, $614$ cells, about $88$ s; the hypothesis is discharged.\\
\bottomrule
\end{tabular}
\end{table}

\subsection{Interval branch-and-bound}\label{ssec:bnb}
\paragraph{Reductions.} Since $k^2\ge0$, $F_{\bm}(g)\ge\sum_i\mu_ig_i$, so a box on which $\sum_i\mu_i(c_i-h_i)$ already exceeds the claim is
discarded, and only a bounded box needs to be covered. The weights and credits of the all-marks certificates are reversal-symmetric, so
$F_{\bm}(g)$ is unchanged when both the pattern and the gap vector are reversed; this reduces the $2^K$ patterns to $20$, $36$ and $72$ classes for
$K=5,6,7$. Coincident points ($g_i=0$) belong to the domain; there $k(0)=1$ makes $F$ large.

\paragraph{Kernel bounds.} For a window $\psi\ge0$, $\abs{k_\psi^{(j)}}\le(2\pi)^j\int v_\psi\abs s^j$; for $\tilde\psi$ this gives
$\abs{k'''}\le M_3=\frac{1747823}{7775424}\pi^3<6.969843$. These are proved analytically and enter as constants. Values of $k,k',k''$ at box centres are enclosed in Arb (or in
\texttt{mpmath.iv} with explicit Taylor remainders).

\paragraph{Lower bounds on a box} $\mathbf c\pm\mathbf h$. For a span $s$ (a run of consecutive gaps) let $\mathrm{sc}_s$, $\mathrm{hs}_s$ be the
sums of the $c_i$ and of the $h_i$ over the span.
\begin{itemize}
\item[(LB1)] On the span interval, $\abs k\ge\abs{k(\mathrm{sc})}-\abs{k'(\mathrm{sc})}\,\mathrm{hs}-\frac12\sup\abs{k''}\,\mathrm{hs}^2$; clip at $0$,
square, sum with the weights, and add $\sum\mu_i(c_i-h_i)$.
\item[(LB2)] $F(\mathbf c)-\sum_i\abs{\partial_iF(\mathbf c)}h_i-\frac12\sum_s\gamma_sW_{2,s}\,\mathrm{hs}_s^2$, where $W_{2,s}$ bounds
$\abs{w''}=2\abs{k'^2+kk''}$ on the span interval; valid because the Hessian of a span term is $w''\,\mathbf v_s\mathbf v_s^{\tsp}$.
\item[(LB3)] By Taylor's theorem with integral remainder,
$F(\mathbf c+\boldsymbol\delta)\ge F(\mathbf c)+\nabla F(\mathbf c)\cdot\boldsymbol\delta+\frac12\sum_s\gamma_s\,\underline w''_s\,(\mathbf v_s\cdot\boldsymbol\delta)^2$
on the box, with $\underline w''_s$ a lower bound for $w''$ on the span interval from a Taylor model of $k,k',k''$ with $\abs{k'''}\le M_3$. Spans with
$\underline w''_s\le0$ are bounded below by $-\frac12\gamma_s\abs{\underline w''_s}\mathrm{hs}_s^2$. The remaining quadratic $q$ is convex, so
$q(\boldsymbol\delta)\ge q(\hat{\mathbf d})+\nabla q(\hat{\mathbf d})\cdot(\boldsymbol\delta-\hat{\mathbf d})$ for \emph{every} $\hat{\mathbf d}$ (a
floating-point quadratic programme chooses $\hat{\mathbf d}$, which affects only tightness), and this affine function is minimised over the box exactly.
\end{itemize}
A box is accepted when a lower bound, evaluated in interval arithmetic, exceeds the claim; otherwise it is bisected. The certificate is the
statement that the stack empties with every box domain-pruned or accepted. LB3 cuts the trees by factors of $30$ to $200$. Tangent planes
after a convexity certificate are used in the verifiers of \mbox{sxuff}, \mbox{AMTOPA} and \mbox{npip99} (\S\ref{ssec:others}); LB3 is our form of this step, and we
have not compared it with theirs in substance.

\paragraph{Arithmetic and trust.} Box edges are exact dyadics. The default engine uses float64 intervals widened outward by one unit in the last
place at every operation, which assumes IEEE-754 round-to-nearest; a scalar engine in Arb balls avoids this assumption at about fifteen times the
cost. The correctness of Arb is trusted. The closing arithmetic ($H$, $\Phi_m$, the constants) was done in Arb from exact rationals and is also
checked in Lean.

\paragraph{Validation and negative controls.} Claims set just above the located minimum must fail, and do, with a rigorous counterexample:
for $K=5$, class $11111$ at the claim raised by $4.5\cdot10^{-5}$; for $K=6$, class $211112$ at $0.0166770$ ($F=0.01667680$); for $K=7$ the
certifier refuted class $1121212$ at the claims first proposed by the linear programme that chose the weights ($F=0.02054932<0.02054941$), in a
basin that the programme's search had missed, after which every claim was lowered by $10^{-5}$ and all $72$ classes certified; at the
lowered claim plus $10^{-5}$ the certifier refutes class $2112112$. For the polynomial window, the same engine at $K=7$ refutes the claim $0.00736$
($F=0.00735942$), and its certificate at the claim $0.00732$ was reproduced box for box independently ($18{,}215{,}893$ boxes). The
$K=8$ claim $0.00796$ lies $0.56\%$ below the floating-point minimum $0.0080050$. An independent kernel stress test of LB3 on $2400$ boxes found no
violation. The $K=8$ and $K=7$ certificates were each re-run in full: independent runs, not independent code.

\subsection{Kernel replay inside Lean}\label{ssec:replay}
\paragraph{The majorant (done).} The checker is written in core Lean. For each cell it proves one enclosure of $(\sin\theta,\cos\theta)$, $\theta=\pi m/30$
at the cell midpoint $m$, and generates the other $58$ cosines by a fixed-point complex recurrence $z_{k+1}=\lfloor z_kz_1/2^{64}\rfloor$ whose error
grows only linearly ($\abs{z_k-e^{ik\theta}}\le k\cdot2^{-55}$, proved); the sums $P(m)$, $P'(m)$, $P''(m)$ are exact integer sums, and a third-order
bound with $\abs{P'''}\le645$ covers the cell. $P\ge0$ on $[0,30]$ needs $600$ cells, $\cB\ge v_\psi$ on $[0,\frac12]$ needs $14$. The replay takes
about $88$ seconds in $13$ files, at most $1.64$ GB each; the smallest margins are $4.1\cdot10^{-6}$ and $4.4\cdot10^{-6}$. Negative controls fail as
they should, in a python mirror ($62$ cells) and in the kernel, which proves that the checker returns \texttt{false} on $[3.25,3.75]$ and on
$[\frac14,\frac12]$ for deliberately wrong data. The analytic part (the Fourier pair $\sinc^2\leftrightarrow$ triangle, the transform of $\cB$, its
support, $\int\cB=\beta_0$, and $2\beta_0+4\abs{\beta_{52}}<\frac{33}{10}$ in exact rationals) is proved in Lean. So the hypothesis \texttt{CertMajV2}
is discharged.

\paragraph{The $K=5$ all-marks certificate (replayed by the kernel).} The branch-and-bound trees above cannot be replayed box by box in the kernel. A second
certificate with large leaves was built for the same data: $92{,}074$ nodes and $46{,}047$ leaves over the $20$ classes, of three kinds --- cap
leaves (the penalty alone exceeds the claim), convex leaves (a tangent plane at a floating-point minimiser plus an exact positive semidefiniteness
check) and linear-programming leaves (five affine minorants per span with rational dual multipliers; $4{,}660$, $2{,}343$ and $39{,}044$ leaves
respectively). Every leaf is accepted by the exact rational mirror of the corrected checker (version~3, described below; $46{,}047$ leaves,
no failure). In a differential test on mutated leaves the mirror was never more permissive than the Lean checker; in any case soundness
rests on the kernel replay, not on the mirror. The checker stores each kernel enclosure (the values
$k,k',k''$ at a rational point, or bounds on $w,w''$ over a cell of width $\frac1{64}$) as a record proved once by kernel evaluation; leaves only
combine stored rationals. Measured in the replay of class $11111$ ($1{,}725$ leaves): about $0.24$ s per point record, $0.025$ s per cell and
$0.74$ s per leaf; with at most $40$ leaves per file the leaf files peak at $3.43$ GB in the one-class pilot. The full replay ($1{,}478$ files: $1{,}475$ shards and
three common files; $46{,}047$ leaves, $89{,}362$ point records; manifest sha256 \texttt{360d9b22469b9dbf449899a2c6d0bd5b}\allowbreak
\texttt{e696380e826f137ef20da71ae96c1609}) was estimated from the pilot at about $15.6$ Lean-hours; it took $20.4$. At the minimiser of class $11111$ the margin is $1.49\cdot10^{-5}$; raising
the claim by $1.5\cdot10^{-5}$, to $8.9\cdot10^{-8}$ above the true minimum, makes exactly one leaf fail in the kernel, the one that contains
the minimiser, as it should.
The soundness theorem of the checker is split into $28$ nodes. An earlier version of one leaf check was unsound: the linear-programming leaf
took a tangent line at a point computed from the floating-point minimiser without checking that this point lies in the span interval, where
alone the second-derivative penalty is valid; a kernel-checked counterexample (one span, a claim of $\frac9{10}$ accepted where the functional
is about $0.446$) shows it. The majorant certificate uses only the interval and enclosure layer of the checker, which is proved.
Version~3 of the checker corrects the defect with one guard, which admits the tangent line at the floating-point point only when that point
lies in the span interval; no other definition used by the checker changed, and the $20$ class certificates pass it unchanged (through the exact mirror).
The whole soundness chain of version~3 is proved with the three standard axioms and is in the library (Appendix~\ref{app:chain}):
\texttt{ZetaS.CertV2.\allowbreak ChainV3.\allowbreak certAM5\_of\_checks} derives \texttt{CertAM5} from the acceptance of the $20$ class certificates by the checker
and the identity of their data with the data of \texttt{CertAM5}.

The replay files use the library's own definitions. Built from an empty build folder, the one-class pilot compiles $56$ files with no
failure in $1{,}976$ s, at a peak of $3.49$ GB. The full replay, against the library at commit \texttt{bb7223b}, whose Lean files the release
changes only in comments (Appendix~\ref{app:layout}), compiled all $1{,}478$ files, with no failure and no retry, in $2$~h~$50$~min of
wall time and $20.4$ Lean-hours of processor time; the lowest free memory was $2.48$ GB and the build products take $7.1$ GB. It ran in two
pools, $273$ light files sixteen at once and $1{,}203$ heavy files six at once, and then the two top files alone ($21$~s and $18$~s, peak
$4.65$ and $4.67$ GB). The kernel checked the $46{,}047$ leaves and $89{,}362$ point records in exact rational arithmetic. The final file applies
\texttt{certAM5\_of\_checks} with the identity of the data by \texttt{rfl} and proves \texttt{ZetaS.CertV2.\allowbreak certAM5\_replayed};
the two theorems of the $K=5$ pair without hypothesis follow (Appendix~\ref{app:replay}). These three theorems
depend only on \texttt{propext}, \texttt{Classical.choice} and \texttt{Quot.sound} and rest on no open declaration (Appendix~\ref{app:census}). The $1{,}478$ source files
match their manifest file by file, and a scan of all of them finds no \texttt{sorry}, \texttt{native\_decide}, \texttt{axiom},
\texttt{unsafe}, \texttt{implemented\_by} or \texttt{extern}. The $1{,}478$ files, byte-identical to those checked, are the
library \texttt{ZetaSReplay} of the repository, and the two theorems are in its module \texttt{ReplayHeadlines}; \texttt{lake build
ZetaSReplay} repeats the replay (\S\ref{ssec:reproduce}).

\paragraph{The $K=7$ and $K=8$ certificates.} Large leaves cut the $K=7$ trees only by a factor of about four: an estimated $6\cdot10^7$ leaves and
$15$ GB of certificate, because the functional has a lattice of near-minimal basins whose number grows geometrically with $K$ (the branch-and-bound
trees grow by a factor of $20$ to $30$ per step in $K$). This is beyond the $32$ GB machine used. These certificates therefore remain displayed
hypotheses, identified by the sha256 of Table~\ref{tab:certs} and verified outside Lean as described in \S\ref{ssec:bnb}.

\subsection{How to reproduce}\label{ssec:reproduce}
\sloppy
{\raggedright Everything is in the repository \repourl, at tag \repotag{} (Appendix~\ref{app:layout}). The Lean libraries are at its root, the replay
files of \S\ref{ssec:replay} are the library \texttt{ZetaSReplay}, and the folder \texttt{supplementary/} holds the data of
Table~\ref{tab:certs}, the certifiers and their logs, and the generator of the replay files, with a guide (\texttt{supplementary/README.md}),
the commands that run them (\texttt{supplementary/VERIFY.md}) and the sha256 of every file
(\texttt{supplementary/CHECKSUMS.sha256}). Commands are run from the root of the repository.\par}
\begin{itemize}\raggedright
\item \emph{The statements and their axioms.} The toolchain is pinned in \texttt{lean-toolchain}, and
\texttt{lake exe cache get} fetches the prebuilt Mathlib. Then
\begin{Verbatim}[fontsize=\small,xleftmargin=1em]
lake build ZetaS
lake build Solution.FollowUpZeta
lake env lean comparator/PrintAxioms/FollowUpZeta.lean
\end{Verbatim}
Every line printed by the last command must read \texttt{\ldots{} depends on axioms: [propext, Classical.choice, Quot.sound]}. This checks
the fourteen statements of record (Appendix~\ref{app:record}), in which \texttt{CertAM5}, \texttt{CertAM7} and \texttt{CertS8} are
hypotheses. The files \texttt{comparator/README\_followup.md} and \texttt{comparator/config-followup.json} (the fourteen theorem names and
the permitted axioms) describe the same check with the comparator tool used by the repository of~\cite{AF}.
\item \emph{The kernel replay of \texttt{CertAM5}} (about $20$ Lean-hours, and several GB of memory per file, \S\ref{ssec:replay}):
\begin{Verbatim}[fontsize=\small,xleftmargin=1em]
lake build ZetaSReplay
lake env lean comparator/PrintAxioms/FollowUpReplay.lean
\end{Verbatim}
The last command prints the axioms of \texttt{ZetaS.Top.zeta\_simple\_K5\_uncond}, \texttt{ZetaS.Top.zeta\_distinct\_K5\_uncond} and
\texttt{ZetaS.CertV2.certAM5\_replayed}; again every line must read as above.
\item \emph{The local certificates outside Lean.} In \texttt{supplementary/certificates/}: \texttt{CertS8/} (the certifier
\texttt{bnb\_lb3.py} with \texttt{polywin.py}; the driver logs of the run and of the re-run, which end with \texttt{ALL SHARDS OK}),
\texttt{CertAM5/} (\texttt{bnb\_interval\_w.py}, the $K=5$ run, with its driver \texttt{runcert.sh}) and \texttt{CertAM7/}
(\texttt{bnb\_lb3\_w.py}, the LB3 engine with the all-marks check, with its driver \texttt{runlb3.sh}, and the exact constants
\texttt{pexact2.py}), each with its data and logs; and \texttt{common/}, with the input checks \texttt{verify\_inputs.py} (weights
nonnegative, span sums $2$, penalties summing to $\nu$, claims exactly $\sum_ib_i(m_i)$, all patterns covered) and \texttt{d819\_exact.py}.
\item \emph{The majorant:} \texttt{supplementary/certificates/CertMajV2/r4\_majorant\_arb.py} (Arb, $106$ bits); in Lean the files
\texttt{MajCheck.lean}, \texttt{MajSoundNPos.lean}, \texttt{MajSoundNMaj.lean} and the $13$ replay shards, all in \texttt{ZetaS/Majorant/}
and built with \texttt{ZetaS}.
\item \emph{The replay files from the certificate:} \texttt{supplementary/replay-K5/} holds the certificate with large leaves
(\texttt{v2\_K5\_*.json}, one file per class, and \texttt{v2common\_K5\_4.json}) and the generator \texttt{gen\_replay.py}, which writes the
replay files from it; \texttt{ZetaSReplay/MANIFEST.sha256} gives the sha256 of each replay file.
\end{itemize}
